\section{Alternative certification methods}\label{ind:section}

The primary certificate does not use the methods in this appendix.
We develop several ways to assess or reproduce its finite quantities.
The first gives a fixed rational evaluation of every $N=256$ moment and
proves a uniform error theorem.  A separate exact-rational
program checks the first $6\times6$ block with an $N=64$ rule.
Rational-center Taylor integration gives a different, arbitrarily refinable
enclosure method, and the final proposition gives a local formula for the
direct cardinal derivatives.

All terminating decimals denote exact rational numbers.  The full
$N=256$ rational recipe and the Taylor method specify computations and
their error bounds without asserting unreported finite values.  The
$N=64$ program has the narrower scope stated with its proof.

\subsection{Rational exponential primitives}\label{cert:rational-primitives}

The fixed rational procedures below use rational arithmetic alone, with a
complex number represented by a pair of rationals.  Put
\[
 \begin{aligned}
 p_*&=3.141592653589793238462643383279502884,\\
 b_*&=.866025403784438646763723170752936183,
 \end{aligned}
\]
and retain $B=4/3$ and $h=2/5$ exactly.  Both $|p_*-\pi|$ and $|b_*-b|$
are less than $10^{-36}$.  For example, the identity
$\pi/4=4\arctan(1/5)-\arctan(1/239)$ and alternating series through
power $57$ certify the first assertion; comparing $(b_*\pm10^{-36})^2$
with $3/4$ certifies the second.
Define
\[
 [w]_{40}=10^{-40}\lfloor10^{40}w\rfloor
\]
for real $w$, and apply it componentwise to complex numbers.  Define an
approximation $e_*(u)$ to $e^{i\pi u}$ by the following finite recursion:
\begin{equation}\label{cert:estar}
 v_0=\left[\sum_{q=0}^{32}\frac{(ip_*u/1024)^q}{q!}\right]_{40},
 \qquad v_{a+1}=[v_a^2]_{40}\ (0\le a<10),\qquad e_*(u)=v_{10}.
\end{equation}
Real and imaginary parts of $e_*(u)$ give the cosine and sine
replacements when $u$ is real.

\begin{lemma}[Rounded exponential]\label{cert:rounded-exponential}
For every $u\in\mathbb Q(i)$ satisfying
$\Re(ip_*u)\le0$ and $|p_*u|\le1000$, the recursion
\eqref{cert:estar} satisfies
\[
 |e_*(u)-e^{ip_*u}|<3.542\cdot10^{-34}<10^{-32}.
\]
\end{lemma}

\begin{proof}
The reduced argument has modulus less than one, so the Taylor remainder
in the first step of~\eqref{cert:estar} is less than $3/33!$.
A componentwise floor contributes less than $2\cdot10^{-40}$ in complex
absolute value.  Each exact exponential subsequently squared has modulus
at most one.  Thus the recurrence
\[
 e_0=3/33!+2\cdot10^{-40},\qquad
 e_{a+1}=2e_a+e_a^2+2\cdot10^{-40}
\]
bounds the accumulated error.  Exact rational evaluation gives
$e_{10}<3.542\cdot10^{-34}$, as claimed.
\end{proof}

\paragraph{Positive exponential upper bounds.}
For $x\ge0$ and an integer $n\ge0$ with $x<n+2$, define
\begin{equation}\label{ind:positive-exponential}
 \mathcal T_n(x)=\frac{x^{n+1}}{(n+1)!}\frac1{1-x/(n+2)},\qquad
 \mathcal E_n(x)=\sum_{k=0}^n\frac{x^k}{k!}+\mathcal T_n(x).
\end{equation}
After the first omitted term, the ratio of each term to its predecessor is
at most $x/(n+2)<1$, so $e^x\le\mathcal E_n(x)$.  These are rational
upper bounds when $x$ is rational.  They will be used both in the fixed
first-block check and in the arbitrary-refinement Taylor method.

\subsection{A fixed rational evaluation}\label{cert:rational}

We now apply the fixed primitives to the full $N=256$ rule, using rational
coefficients and rational $C,m,y$ in the finite domain of
Lemma~\ref{cert:moment-error}.  After defining the rounded sums, we prove
their uniform approximation error.  No matrix, residual, or Bernstein
values are asserted by that error theorem.

In~\eqref{cert:rule}, replace each cosine $\cos(\pi u)$ by $\Re e_*(u)$,
with $u=\alpha_v$ or $u=2a\alpha_v$ as appropriate, and then apply
$[\ ]_{40}$ to the final weight $o_v$ only.  Denote the results by
$o_v^{\#},\tau_{jv}^{\#}$.  In the displayed coefficient list for $P_l$,
replace $b$ by $b_*$; call the resulting numbers $P_l^{\#}$.
Evaluate~\eqref{cert:density} by replacing $\pi$ by $p_*$, $P_l$ by
$P_l^{\#}$, and every exponential by $e_*$, at the approximate point
$t=\tau_{jv}^{\#}$.  Call the resulting density $f_j^{\#}(t)$.
For either these points or the exact atom points $t=t_j$, use
\begin{equation}\label{cert:lambda-z-star}
 \lambda_*=[i/(b_*(t+ih))]_{40},\qquad
 z_*=[-B/(t+ih)-ih]_{40},
\end{equation}
and set
\begin{equation}\label{cert:rational-factors}
 E_l(t)=e_*(tm)\frac{(ip_*ty)^l}{l!},\qquad
 G_l(t)=\lambda_*e_*(z_*m)\frac{(ip_*z_*y)^l}{l!}.
\end{equation}
The rational moment sums are therefore
\begin{equation}\label{cert:rational-moment-sums}
\begin{aligned}
 U_{p,l}^{\#}
 &=\Re\left\{\frac16\sum_{j=1}^6\sum_{v=0}^{255}
    o_v^{\#}f_j^{\#}(\tau_{jv}^{\#})E_l(\tau_{jv}^{\#})
       +C\sum_{j=0}^6e_jP_j^{\#}E_l(t_j)\right\},\\
 X_{p,l}^{\#}
 &=\Re\left\{\frac16\sum_{j=1}^6\sum_{v=0}^{255}
    o_v^{\#}f_j^{\#}(\tau_{jv}^{\#})G_l(\tau_{jv}^{\#})
       +C\sum_{j=0}^6e_jP_j^{\#}G_l(t_j)\right\}.
\end{aligned}
\end{equation}
No additional floors are used.  For $Q_n^{\#},D_n^{\#}$ use the product
and cotangent formulas defining $Q_n,D_n$, with the same sine and cosine
replacements and residue $a=n\bmod12\in A$; cotangent means the ratio of
the replaced cosine to the replaced sine.


For the finite pseudocode in Section~\ref{cert:evaluation}, the rational
prescription uses the same measure order: the $6\cdot256$ quadrature pairs
followed by the seven atom pairs.  Put $\#$ on the prescribed quantities
and make exactly the replacements above; no other rounding is performed.

\begin{lemma}[Rational moment and node errors]\label{cert:rational-moment-error}
Let $p=PR$ have a rational coefficient list supported in
$\{0\}\cup I_f$, with $\|(c,d)\|_1\le3$, and let $C\in\mathbb Q$
satisfy $|C|\le1$.  Let $m,y\in\mathbb Q$ and
$l\in\{0,\ldots,30\}$ satisfy $0\le m\le100$ and $|y|\le3/2$.
The prescription in Section~\ref{cert:rational} gives rational numbers
$U_{p,l}^{\#}$ and $X_{p,l}^{\#}$ such that
\[
 |U_{p,l}^{\#}-U_{p,l}|<10^{-14},\qquad
 |X_{p,l}^{\#}-X_{p,l}|<10^{-14}.
\]
Separately, for every node $n\in12\ZZ+A$, using its residue $a\in A$
in~\eqref{fourier:node-data}, the prescribed rational values satisfy
\[
 |Q_n^{\#}-Q_n|<10^{-25},\qquad |D_n^{\#}-D_n|<10^{-25}.
\]
\end{lemma}

\begin{proof}
We verify the call hypotheses of Lemma~\ref{cert:rounded-exponential}
before using its error bound, and then propagate the errors through the
finite sums.  The cosine calls defining the nodes and weights have real arguments
$u=\alpha_v$ or $u=2a\alpha_v$, with $|u|<254$; hence
$|p_*u|<3.142\cdot254<1000$ and $\Re(ip_*u)=0$.
For real $u$, replacing $p_*$ by $\pi$ adds at most $|u|10^{-36}$.
The telescoping coefficient sum in the weights is less than one, so
\[
\begin{aligned}
 |\tau_{jv}^{\#}-\tau_{jv}|
  &<\frac{3.542\cdot10^{-34}+10^{-36}}{12}<9\cdot10^{-34},\\
 |o_v^{\#}-o_v|
  &<\frac{3.542\cdot10^{-34}+254\cdot10^{-36}}{256}
       +10^{-40}<\frac{10^{-29}}{256}.
\end{aligned}
\]
The exact nodes have a uniform distance from their piece endpoints:
\[
\begin{aligned}
 \min\{\tau_{jv}-t_{j-1},t_j-\tau_{jv}\}
 &\ge\frac{1-\cos(\pi/512)}{12}
  =\frac{\sin^2(\pi/1024)}6\\
 &\ge\frac1{6\cdot512^2}=\frac1{1572864}>9\cdot10^{-34}.
\end{aligned}
\]
Here $\sin x\ge2x/\pi$ on $[0,\pi/2]$ follows from concavity.
It follows that every real rounded node $\tau_{jv}^{\#}$ remains strictly
inside $(t_{j-1},t_j)$.  The density is therefore evaluated on its stated
piece, and all rounded nodes and exact atoms lie in $[0,1]$.
The remaining real exponential calls in the density and direct factors
now have $|u|\le100$; the residue arguments for $Q_n,D_n$ have $|u|<1$.
They too satisfy the call hypotheses.

For real $t\in[0,1]$, the exact identities
\[
\begin{gathered}
 |\lambda(t)|\le\frac5{\sqrt3}<3,\\
 |z(t)|^2=h^2+\frac{B^2-2Bh^2}{t^2+h^2}\le\left(\frac{44}{15}\right)^2,\\
 \operatorname{Im}z(t)=h\left(\frac B{t^2+h^2}-1\right)
       \ge\frac{26}{435}>.0597
\end{gathered}
\]
use $B^2-2Bh^2=304/225>0$.  They apply in particular at every rounded
node before the floor in~\eqref{cert:lambda-z-star}.  That floor lowers the
imaginary part by less than $10^{-40}$ and changes the complex value by
less than $2\cdot10^{-40}$.  Thus
\[
 \operatorname{Im}z_*\ge\frac{26}{435}-10^{-40}>0,
 \qquad |z_*|<\frac{44}{15}+2\cdot10^{-40}<3.
\]
For the exact transformed arguments, $100\pi|z(t)|<922$; for the rounded
ones, $|p_*z_*m|<100(3.142)(3)<1000$ and
$\Re(ip_*z_*m)=-p_*m\operatorname{Im}z_*\le0$.
The exponential estimate already proved therefore applies to every
transformed call as well.

For a quadrature node write $t=\tau_{jv}$ and $t^{\#}=\tau_{jv}^{\#}$;
at an atom put $t=t^{\#}=t_j$.  Along the intervening real segment,
$|t+ih|\ge h$ and $b,b_*>.86$.  The derivative bounds for the reciprocal
factors and the change of $b$ give
\[
\begin{aligned}
 |\lambda_* -\lambda(t)|
  &<8(9\cdot10^{-34})+4\cdot10^{-36}+2\cdot10^{-40}<10^{-29},\\
 |z_*-z(t)|
  &<\frac{25}{3}(9\cdot10^{-34})+2\cdot10^{-40}<10^{-29}.
\end{aligned}
\]
Indeed $1/(b_*h^2)<8$, $B/h^2=25/3$, and the change of $1/b$ contributes
less than $10^{-36}/(.86^2h)<4\cdot10^{-36}$.  At atoms the node-error
terms vanish.  Both transformed exponential arguments lie in the closed
left half-plane after multiplication by $i$.  The integral first-difference
formula there gives, for example,
\[
 |e_*(z_*m)-e^{i\pi z(t)m}|
 <10^{-32}+100(3.142\cdot10^{-29}+3\cdot10^{-36})<4\cdot10^{-27}.
\]
The direct exponential-factor error is smaller.  The arguments of the
divided powers differ by less than
$1.5(3.142\cdot10^{-29}+3\cdot10^{-36})<5\cdot10^{-29}$.
The divided powers and their derivatives are bounded along the connecting
segment by $e^{1.5\cdot3.142\cdot3}<1.5\cdot10^6$; this also gives the
stated bound $e^{1.5\pi\cdot3}<1.5\cdot10^6$.  One rational check of this
exponential bound uses $1.5(3.142)(3)/16<.884$ and
\[
\begin{gathered}
 \sum_{q=0}^{8}\frac{.884^q}{q!}
 +\frac{.884^9}{9!}\frac1{1-.884/10}<2.421,\\
 2.421^{16}<1.5\cdot10^6.
\end{gathered}
\]
The first expression majorizes the exponential series because every
successive ratio in its tail is at most $.884/10$.
The divided-power errors are therefore less than
$7.5\cdot10^{-23}<1.5\cdot10^{-22}$.
Using $|\lambda|<3$, modulus at most one for the exact exponentials, and
modulus less than two for their approximations, the transformed-factor
error is at most
\[
 2(1.5\cdot10^6)10^{-29}
 +3(1.5\cdot10^6)(4\cdot10^{-27})
 +3(1.5\cdot10^{-22})<3\cdot10^{-20}.
\]
The same bound applies to the direct factor.

For completeness, the density estimate can be obtained without relying on
cancellation.  Its real exponential factors, including evaluation at a
rounded node, have error less than
$3.542\cdot10^{-34}+100(3.142(9\cdot10^{-34})+10^{-36})<3\cdot10^{-31}$.
Let $A_n,B_n$ denote the inner sums in~\eqref{cert:density} with and without
the factor $t_l-t$, respectively, and put $\#$ on their replacements.
On the common piece, $|t_l-t|,|t_l-t^{\#}|\le1$ for $l\ge j$.
The displayed list for $P_l$ gives
$\sum_{l=1}^6|P_l^{\#}-P_l|<7\cdot10^{-36}$ and
$\sum_{l=1}^6|P_l^{\#}|<11$.  Hence $|A_n^{\#}|,|B_n^{\#}|<22$ and
\[
\begin{aligned}
 |A_n^{\#}-A_n|,\ |B_n^{\#}-B_n|
 &<22(9\cdot10^{-34})+2(7\cdot10^{-36})\\
 &\quad+10(3\cdot10^{-31})<4\cdot10^{-30}.
\end{aligned}
\]
The bracket in~\eqref{cert:density} has approximate magnitude at most
$70(|c_n|+|d_n|)$ and error less than
$1.3\cdot10^{-29}(|c_n|+|d_n|)$.  Expanding the outer product once and
using list norm at most three gives the pointwise bound
\[
\begin{aligned}
 |f_j^{\#}(t^{\#})-f_j(t)|
 &<2(10^{-36})(3)(2)(70)\\
 &\quad+2(3.142)(3)\bigl(70(3\cdot10^{-31})+1.3\cdot10^{-29}\bigr)\\
 &<7\cdot10^{-28}<10^{-22}.
\end{aligned}
\]

On the real segment $|f_j|<1000$, and each atom has magnitude at most six
and error less than $10^{-28}$.  The latter error also follows directly
from the displayed $P_l$ list.  The exact transformed factor has magnitude
below $3\cdot1.5\cdot10^6$, so every approximate factor has magnitude
below $4.6\cdot10^6$.  The exact quadrature weights have total mass one,
and the sum of the absolute weight errors over all six pieces is less than
$10^{-29}$.  Separating weight, density, factor, and atom errors, and then
adding the analytic quadrature error, bounds either moment error by
\[
\begin{split}
 &10^{-29}(1001)(4.6\cdot10^6)+10^{-22}(4.6\cdot10^6)
   +1000(3\cdot10^{-20})\\
 &\quad+7\bigl(10^{-28}(4.6\cdot10^6)+6(3\cdot10^{-20})\bigr)
   +10^{-22}<5\cdot10^{-16}<10^{-14}.
\end{split}
\]

Finally, for the residue arguments the sine and cosine replacement errors
are less than $3.542\cdot10^{-34}+10^{-36}<4\cdot10^{-34}$.
Every nonzero exact sine has magnitude at least $\sin(\pi/12)>1/4$, and
its replacement has magnitude greater than $1/5$.  Each squared factor
$(2\sin)^2$ in $Q_n$ and its replacement is less than five in magnitude,
and their difference is less than $5\cdot10^{-33}$.  Telescoping the five
factors and allowing for $(p_*/6)^2-(\pi/6)^2$ gives a $Q_n$ error less
than $2\cdot10^{-29}<10^{-25}$.  For a cotangent, division by the two sine
lower bounds gives an error less than
$5(4\cdot10^{-34})+20(4\cdot10^{-34})=10^{-32}$.
Summing five terms and allowing for the prefactor gives a $D_n$ error less
than $6\cdot10^{-32}<10^{-25}$, as required.
\end{proof}

\subsection{Propagation to finite certificates}\label{cert:propagation}

We record stability estimates before applying any numerical margin.
The first controls the two operator coordinates under errors in the moments
and node constants.  The second controls Bernstein sums from errors in
the signed-$y$ moments.

To propagate a moment radius, abbreviate $X_\ell=X_{p,\ell}(m,1)$ and
write $\mathcal S_{m,c},\mathcal S_{m,d}$ for the two components of
\eqref{cert:pair}, with hats for their rational approximations.
Suppose $|X_0|,|X_1|\le H$,
$|D_m|\le d$, and rational approximations have errors at most
$\epsilon,\epsilon,\eta_Q,\eta_D$ in $X_0,X_1,Q_m,D_m$, respectively.
If both exact and approximate $Q_m$ exceed $q>0$, expanding the numerator
and then the reciprocal in \eqref{cert:pair} gives
\begin{equation}\label{ind:pair-stability}
 \max\{|\widehat{\mathcal S}_{m,c}-\mathcal S_{m,c}|,
        |\widehat{\mathcal S}_{m,d}-\mathcal S_{m,d}|\}
 \le \frac{(d+1+\eta_D)\epsilon+H\eta_D}{q}
       +\frac{(d+1)H\eta_Q}{q^2}.
\end{equation}
For example, the perturbed second numerator differs by at most
$(d+\eta_D)\epsilon+H\eta_D+\epsilon$, and the exact numerator is bounded
by $(d+1)H$; this proves the inequality, which also dominates the first
entry's error.  The Taylor budgets and the separated and common-kernel
first-block estimates below use $q=7/4$ and $d=2.12$.  The separate-rounding
comparison uses $q=7/4$ and $d=4.3$.
The rational lower bound $Q_m>197192/112500$ proved after
\eqref{fourier:node-bounds} exceeds $7/4$ by $1268/450000$, so
$\eta_Q\le10^{-25}$ also ensures that the approximate $Q_m$ exceeds $q$.

\paragraph{Analytic pair-error allowance.}
For an input $p$ in the rational finite-input domain of
Lemma~\ref{cert:rational-moment-error} and a row $m\in I_f$, define the rational
version $\mathcal S_m^\#$ by using the quantities with superscript $\#$
throughout~\eqref{cert:pair}.  We record the division estimate explicitly
for both basis and full-table inputs.  By
\eqref{interpolation:piece-envelopes} and the atom masses there, the total
variation of the folded measure for any input in the real finite-input
domain of Lemma~\ref{cert:moment-error} is at most
\[
 3\sum_{j=1}^6\frac{M_j}{6}+\sum_{j=0}^6m_j^*=119.642<133.
\]
On the real segment, $|\lambda|,|z|<3$ and
$|e^{i\pi z m}|\le1$ for $m\ge0$.  Thus, using $\pi<22/7$,
\[
 |X_{p,0}(m,1)|<133(3)<400,\qquad
 |X_{p,1}(m,1)|<133(3)(3)(22/7)=3762<4000.
\]
For this calculation abbreviate $X_j=X_{p,j}(m,1)$, $Q=Q_m$, and $D=D_m$.
Write $\varepsilon=10^{-14}$ and $\delta=10^{-25}$.  By
Lemma~\ref{cert:rational-moment-error}, the errors in $X_0,X_1$ are less than
$\varepsilon$ and those in $Q,D$ are less than $\delta$.
Equation~\eqref{fourier:node-bounds} gives $Q>1.75$, $|D|<2.12$, and hence
$Q^{\#}>1.7$.  For $\mathcal A=DX_0-X_1$ and
$\mathcal A^{\#}=D^{\#}X_0^{\#}-X_1^{\#}$,
\[
 |\mathcal A|<4848,\qquad
 |\mathcal A^{\#}-\mathcal A|<2.12\varepsilon+(400+\varepsilon)\delta+\varepsilon
             <3.121\cdot10^{-14}.
\]
The two quotient differences are consequently bounded by
\[
\begin{aligned}
 \left|\frac{X_0^{\#}}{Q^{\#}}-\frac{X_0}{Q}\right|
 &<\frac{10^{-14}}{1.7}+\frac{400\cdot10^{-25}}{1.75(1.7)}
   <6\cdot10^{-15},\\
 \left|\frac{\mathcal A^{\#}}{Q^{\#}}-\frac{\mathcal A}{Q}\right|
 &<\frac{3.121\cdot10^{-14}}{1.7}
   +\frac{4848\cdot10^{-25}}{1.75(1.7)}
   <1.84\cdot10^{-14}<2\cdot10^{-14}.
\end{aligned}
\]
In particular, the uniform error in each entry of $\mathcal S_m^{\#}$ is
strictly less than $10^{-12}$; we retain this looser allowance below.

\subsubsection{Bernstein multipliers}\label{ind:bernstein}
For every signed half-gap in Lemma~\ref{signs:bernstein}, $y\ne0$ and
$|y|^{-\nu}\le4$.  Directly from \eqref{cert:moments},
\[
 y^{-\nu}\bigl(U_{p,r+\nu}(m,y)+X_{q,r+\nu}(m,y)\bigr)
 =y^r\bigl(U_{p,r+\nu}(m,1)+X_{q,r+\nu}(m,1)\bigr).
\]
This identity also holds for negative $y$ and gives exactly the
power coefficients of \eqref{cert:power-from-moments}; it does not replace
their derivatives by interpolated jets.  Two moment errors less than
$10^{-14}$ give a $T_r$ error less than $8\cdot10^{-14}$.
For fixed $0\le k\le28$, the rational recurrence $a_0=1$ and
$a_r=a_{r-1}(k+1-r)/(29-r)$ satisfies
\begin{equation}\label{ind:bernstein-multiplier}
 a_r=\prod_{j=1}^r\frac{k+1-j}{29-j}
 =\frac{k!}{(k-r)!}\frac{(28-r)!}{28!}
 =\frac{\binom{k}{r}}{\binom{28}{r}}\qquad(0\le r\le k).
\end{equation}
Each factor is in $[0,1]$.  There are
$\sum_{k=0}^{28}k=28\cdot29/2=406$ nonconstant multipliers, in addition
to $a_0=1$ for each $k$.  Every Bernstein sum has at most $29$ terms,
so its error is strictly less than
$29(8\cdot10^{-14})=2.32\cdot10^{-12}$.

\subsection{An independent rational check of the first block}
\label{cert:rational-small}

The optional program \path{verification/check_sjj_rational.py} checks
the $6\times6$ block on $J=\{1,3,4\}$ by exact rational arithmetic.
It uses the cosine rule \eqref{cert:rule} with $N=64$ and the sum over
$1\le a\le31$, and otherwise the rounded exponential, kernels and
basis densities specified above. There are no atom terms for these
basis inputs. The following error bound makes this smaller check
independent of the full Arb calculation.

\begin{lemma}\label{cert:rational-small-error}
For each basis input $c_n=1$ or $d_n=1$, $n\in J$, with all other data
zero, and each output node $m\in J$, the rational program computes
both entries of $\mathcal S_m(p)$ with absolute error less than
$10^{-12}$.
\end{lemma}

\begin{proof}
We separate quadrature error from rounding in the finite sum. Use the
same complex disks of radius $1/6$ as in
Lemma~\ref{cert:moment-error}. The basis density on any disk is bounded by
\[
 D_0=2\pi^2\cdot10\cdot\frac{13}{12}e^{4\pi/6}.
\]
This also bounds a simple-pole basis density. The bounds
$|\lambda|<5$, $|z|<6.2$ and $\operatorname{Im}z>-.081$ proved there
show that the transformed value and first-derivative integrands are
bounded respectively by
\[
 M_0=5D_0e^{4\pi\cdot.081}<24028,
 \qquad M_1=6.2\pi M_0<468015.
\]
Use the rational upper bound $e^x\le\mathcal E_n(x)$ from
\eqref{ind:positive-exponential}.  Put
$\overline\pi=p_*+10^{-36}>\pi$.  Here
$x=743\overline\pi/750$ lies in $[0,4)$, so the formula applies with $n=40$.
Using $M_0=(325/3)\pi^2e^{743\pi/750}$, the exact rational comparisons
\[
 \frac{325}{3}\overline\pi^2 \mathcal E_{40}(743\overline\pi/750)<24028,
 \qquad
 \frac{31}{5}\overline\pi\,
 \frac{325}{3}\overline\pi^2 \mathcal E_{40}(743\overline\pi/750)<468015
\]
prove the displayed bounds for $M_0,M_1$.  Each comparison is a fraction
inequality after clearing positive denominators; no rounded decimal
evaluation of an exponential is used.
Discrete cosine orthogonality now gives exactness through degree $63$.
The weights sum to one and each is at least $1/(64\cdot63)>0$ by
the same telescoping identity. Cauchy's estimate on the half-radius
real segments bounds the degree-$63$ Taylor tails by $M_i2^{-63}$.
The integral and quadrature have total mass one, after combining all
six pieces. Thus their discrepancies in the two moments are at most
\[
 \epsilon_0=2^{-62}M_0<5.211\cdot10^{-15},\qquad
 \epsilon_1=2^{-62}M_1<1.015\cdot10^{-13}.
\]
Using $Q_m>1.75$ and $|D_m|<2.12$, the quadrature error in either
entry of \eqref{cert:pair} is consequently less than
\begin{equation}\label{cert:rational-small-quadrature}
 \frac{2.12\epsilon_0+\epsilon_1}{1.75}<6.432\cdot10^{-14}.
\end{equation}

For completeness, two changes in the rational implementation require
care in reusing the rounding estimates. First, for a real argument $u$
the program evaluates $e_*(r)$ with $r=u\bmod2\in[0,2)$.
Periodicity is used for the exact target:
$e^{i\pi r}=e^{i\pi u}$. The approximation $p_*$ is not treated as
giving an exactly periodic exponential.  Since $r$ is real and
$|p_*r|<2(3.142)<1000$, Lemma~\ref{cert:rounded-exponential}, followed by
$|e^{ip_*r}-e^{i\pi r}|\le2\cdot10^{-36}$, gives an absolute error
less than $2\cdot10^{-32}$ for each such replacement. When $u$
contains an approximate quadrature node, first compare at that
approximate argument and then use the first-difference bound for the
exact exponential. This avoids any assumption of continuity of the
modulo operation.

Second, with $N=64$, the sum of the absolute cosine coefficients in
one weight is less than one. Including the final floor, the weight
error is therefore at most
\[
 \frac{2\cdot10^{-32}}N+10^{-40}
 <\frac{10^{-29}}N.
\]
Summing the stronger first expression, rather than the relaxed bound on
the right, gives the total weight error
\begin{equation}\label{cert:rational-small-weight-total}
 N\left(\frac{2\cdot10^{-32}}N+10^{-40}\right)
 =2\cdot10^{-32}+64\cdot10^{-40}<3\cdot10^{-32}.
\end{equation}
The node error is less than $2\cdot10^{-32}/12<2\cdot10^{-33}$.
The exact nodes are at distance at least $1/98304$ from the endpoints
of their own interval $[(j-1)/6,j/6]$, since
$1-\cos(\pi/128)\ge2/128^2$. Thus the approximate nodes remain in their
own intervals, in particular in $[0,1]$. Their transformed frequencies
have positive imaginary part even after the floor, because
$\operatorname{Im}z\ge26/435$ there.  The real-segment estimates in the
proof of Lemma~\ref{cert:rational-moment-error} also give $|z_*|<3$.
Thus every transformed call has its argument in the closed left
half-plane and satisfies $|p_*z_*m|<3.142(3)(4)<1000$.

Here is a direct propagation bound for these smaller inputs. The
real-segment estimates are $|\lambda|<3$, $|z|<3$, and basis-density
modulus less than $214$. For one piece put
\[
 A_n=\sum_{l=j}^6P_le^{i\pi t_ln},\qquad
 B_n=\sum_{l=j}^6t_lP_le^{i\pi t_ln}.
\]
Their moduli are at most ten. Replacing $b$ and the real exponentials
changes each sum by less than $3\cdot10^{-31}$; including the node
movement changes $B_n-tA_n$ by less than $7\cdot10^{-31}$.
The phase $e^{-i\pi tn}$ changes by less than $5\cdot10^{-32}$,
since $n\le4$. Substitution in \eqref{fourier:folded} therefore gives
a density error less than $3\cdot10^{-29}$.
The derivatives of the rational kernels, using $b,b_*>4/5$ and
$|t+ih|\ge2/5$, give errors less than $2\cdot10^{-32}$ in both
$\lambda$ and $z$, including their floors. The first-difference
formula for exponentials in the closed left half-plane, with $m\le4$,
then bounds the errors in
\[
 g_0(t)=\lambda(t)e^{i\pi mz(t)},\qquad
 g_1(t)=g_0(t)i\pi z(t)
\]
by $2\cdot10^{-30}$ and $3\cdot10^{-29}$, respectively.
The exact factors have moduli below $3$ and $36$; their approximations
have moduli below $4$ and $37$. Product differences bound each
derivative-integrand error by
\[
 3\cdot10^{-29}\cdot37+214\cdot3\cdot10^{-29}
 <8\cdot10^{-27};
\]
the value-integrand bound is smaller. By
\eqref{cert:rational-small-weight-total}, the total absolute weight error
is below $3\cdot10^{-32}$, and the approximate integrands have modulus
below $10^4$. Hence each rounded moment differs from its exact finite
quadrature moment by less than $10^{-26}$.

Finally, for residues $1,3,4$, the exact polynomial forms are
\[
\begin{gathered}
 Q_1=Q_3=\frac{4\pi^2}{3}(1-b),\qquad Q_4=\frac{\pi^2}{3},\\
 D_1=\frac{\pi(3+2b)}{18},\qquad
 D_3=-\frac{\pi(3+2b)}{18},\qquad D_4=\frac{7\pi b}{9}.
\end{gathered}
\]
The program replaces $\pi,b$ by $p_*,b_*$ in these expressions.
Using $|p_*^2-\pi^2|<6.3\cdot10^{-36}$, $\pi^2<10$, and $b,b_*<1$ gives
\[
\begin{aligned}
 |Q_m^\#-Q_m|
 &<\max\{(4/3)(6.3+10),\,6.3/3\}10^{-36}<3\cdot10^{-35},\\
 |D_m^\#-D_m|
 &<\max\{(5+44/7)/18,\,29/9\}10^{-36}<4\cdot10^{-36}.
\end{aligned}
\]
The exact rational lower bound used
after~\eqref{ind:pair-stability} therefore shows that both denominators
$Q_m,Q_m^\#$ exceed $7/4$. The moments have modulus below $10^4$, so the
additional division and $D_m$ errors are less than $4\cdot10^{-31}$.
Together with \eqref{cert:rational-small-quadrature}, these bounds
give an error less than $7\cdot10^{-14}<10^{-12}$ in either matrix
entry, as asserted.
\end{proof}

If every entry of this $6\times6$ matrix has error below
$\varepsilon=10^{-12}$, its induced column-norm error is below
$6\varepsilon$. The program accordingly enlarges the computed
defect norm by $\|W_\eta\|_1\,6\varepsilon$, then verifies that the
result is below $.003$ and that the resulting inverse bound is below
$32$. This is the same perturbation-of-identity argument as in
Section~\ref{cert:matrix-bounds}. It checks both small inverse
certificates without asserting that the full $N=256$ rational
prescription has been executed.


\subsection{Comparison of first-block majorants}\label{ind:small}

The preceding $N=64$ proof separates the value and derivative integrands.
Two coarser arguments instead use one majorant for both.  They concern the
same basis inputs $n\in J$, row nodes $m\in J$, $y=1$, and orders zero and
one, with no atoms.  The table compares their disk and node bounds with
the resulting entry errors:
\begingroup
\small
\[
\begin{array}{c|c|c|c}
 \text{route}&\text{disk majorant}&|D_m|&\text{entry error}\\ \hline
 \text{separated}&M_0<24028,\ M_1<468015&<2.12&<7\cdot10^{-14}\\
 \text{common kernel}&M=10^6&<2.12&<4.1\cdot10^{-13}\\
 \text{separate rounding}&M=750000&<4.3&<5.23\cdot10^{-13}
\end{array}
\]
\endgroup
The separated route uses the explicit density bound $D_0$ from
Section~\ref{cert:rational-small}.  The two common routes use density and
exponential bounds $(1800,2.8)$ and $(2500,3)$, respectively.

\subsubsection{A common kernel majorant}\label{ind:small-kernel}

The rational comparisons
\begin{equation}\label{ind:small-exponentials}
 \mathcal E_{32}(44/21)<42/5,\qquad
 \mathcal E_{32}(891/875)<14/5
\end{equation}
follow by clearing positive denominators in~\eqref{ind:positive-exponential};
both arguments are nonnegative and less than $34$.  Since
$4\pi/6<44/21$ and $4\pi(.081)<891/875$, the two basis densities on the
complex disks satisfy
\[
\begin{aligned}
 |u_n(t)|&<2(22/7)^2(65/6)(42/5)=12584/7<1800,\\
 |v_n(t)|&<2(22/7)10(42/5)=3696/7<1800.
\end{aligned}
\]
The bounds $|\lambda|<5$, $|e^{i\pi zm}|<2.8$, and
$\pi|z|<6.2\pi<20$ therefore bound both transformed integrands by
$1800(5)(2.8)(20)=504000<10^6$.  The positive degree-63 rule and the
half-radius Cauchy estimate give, for each moment,
\begin{equation}\label{ind:small-kernel-quadrature}
 |X_\ell-X_{\ell,\mathrm{quad}}|\le10^6 2^{-62}<2.17\cdot10^{-13}.
\end{equation}
The factor two here compares the integral and quadrature averages; a direct
Taylor average has only one tail, as in Lemma~\ref{ind:taylor-average}.

The rounding proof in Section~\ref{cert:rational-small} applies to these
same inputs and gives
finite moment-rounding error below $10^{-26}$, hence below $10^{-14}$.
It also gives $|X_0|,|X_1|<10^4$, $|\delta Q|<3\cdot10^{-35}$,
$|\delta D|<4\cdot10^{-36}$, and both denominators above $7/4$.
Thus the total moment allowance is $2.27\cdot10^{-13}$.  In
\eqref{ind:pair-stability}, the additional $Q,D$ contribution is less than
\[
 \frac{10^4(4\cdot10^{-36})+(4\cdot10^{-36})(2.27\cdot10^{-13})}{7/4}
 +\frac{3.12\cdot10^4(3\cdot10^{-35})}{(7/4)^2}<4\cdot10^{-31}.
\]
Consequently
\begin{equation}\label{ind:small-kernel-entry}
 \frac{3.12(2.27\cdot10^{-13})}{1.75}+4\cdot10^{-31}
 <4.1\cdot10^{-13}<10^{-12}.
\end{equation}

\subsubsection{A separate rounding majorant}\label{ind:small-rounding}

The weaker bounds $e^{4\pi/6}<9$ and $e^{4\pi(.081)}<3$ also follow
from~\eqref{ind:small-exponentials}.  They bound either basis density by
$2500$; for the larger density,
$2(22/7)^2(65/6)9<2500$.  With the kernel bound $5$ and derivative
factor $6.2\pi<20$, this gives the distinct common majorant
$M=2500(5)(3)(20)=750000$.  Its moment quadrature radius is
\begin{equation}\label{ind:small-rounding-quadrature}
 \varepsilon_q=750000\,2^{-62}
 <1.626303258728257\cdot10^{-13}.
\end{equation}
Using $|D_m|<4.3$ and $Q_m>7/4$, the pair quadrature error is less than
$(5.3/(7/4))\varepsilon_q<4.926\cdot10^{-13}$.
The same finite rounding allowance below $10^{-14}$ applies because the
schedule and inputs have not changed.  Substitution of the same polynomial
$Q,D$ errors in~\eqref{ind:pair-stability}, now with $d=4.3$, bounds their
additional contribution by $6\cdot10^{-31}$.  Therefore
\begin{equation}\label{ind:small-rounding-entry}
 \frac{5.3}{7/4}\bigl(750000\,2^{-62}+10^{-14}\bigr)
       +6\cdot10^{-31}<5.23\cdot10^{-13}<10^{-12}.
\end{equation}

All three routes support the program's $10^{-12}$ entry radius.  They do
not extend its coverage: the public rational algorithm still checks only
the 36 first-block entries and the two inverse certificates, not the
residual or Bernstein tables.

\subsection{Rational-center Taylor integration}\label{ind:taylor}

For a finite input in the real moment domain of Lemma~\ref{cert:moment-error}, let
$F^U_{j,\ell}$ and $F^X_{j,\ell}$ be the density \eqref{cert:density}
multiplied by the respective factors in \eqref{cert:analytic-factors}.
The input is supported on nodes at most $100$, has coefficient-list norm at
most $3$ and $|C|\le1$, and the parameters satisfy
$0\le m\le100$, $|y|\le3/2$, and $\ell\in\{0,\ldots,30\}$.  The coefficients,
$C,m,y$ are rational when a rational enclosure is requested.  The six
centers and the two radii are
\begin{equation}\label{ind:centers}
 c_j=\frac{2j-1}{12},\qquad R=\frac16,\qquad a=\frac1{12}
 \quad(1\le j\le6).
\end{equation}
Each $F^U_{j,\ell}$ and $F^X_{j,\ell}$ is analytic on a neighborhood of
$|t-c_j|\le R$: the input list is finite and the only rational-kernel pole,
$-ih$, is outside that disk.  This analytic assertion is made for finite
inputs.  An arbitrary $\ell^1$ list need not give a complex-frequency
density analytic on these disks.

\begin{lemma}[Taylor averages with rational coefficient radii]
\label{ind:taylor-average}
Suppose $F(c_j+x)=\sum_{k\ge0}A_kx^k$ is analytic on a neighborhood of
$|x|\le R$ and $|F|\le M$ there.  Let $D\ge0$ be an integer.  For
$0\le k\le D$, suppose
$\widehat A_k\in\mathbb Q(i)$ and $e_k\in\mathbb Q_{\ge0}$ satisfy
$|A_k-\widehat A_k|\le e_k$.  Then
\begin{align}\label{ind:taylor-average-error}
 \left|6\int_{t_{j-1}}^{t_j}F(t)\,\dd t-
       \sum_{2r\le D}\widehat A_{2r}\frac{a^{2r}}{2r+1}\right|
 \le M2^{-D}+\sum_{2r\le D}e_{2r}\frac{a^{2r}}{2r+1}.
\end{align}
If the final sum of radii is at most $\delta$ for each of the six pieces,
their combined integral has error at most $M2^{-D}+\delta$.  Discrete
atoms are not included in this estimate and are to be enclosed separately.
\end{lemma}

\begin{proof}
Since $2a=1/6$, the expression $6\int F$ is the interval average.  The
average of $x^k$ on $[-a,a]$ is zero for odd $k$ and is $a^k/(k+1)$ for
even $k$.  Thus the degree-$D$ Taylor polynomial has the exact average
\begin{equation}\label{ind:even-average}
 \sum_{2r\le D} A_{2r}\frac{a^{2r}}{2r+1}.
\end{equation}
Cauchy's estimate gives $|A_k|\le M R^{-k}$.  Bounding the average of the
absolute tail by its supremum on $[-a,a]$ gives the conservative bound
\[
 M\sum_{k>D}(a/R)^k=M\sum_{k>D}2^{-k}=M2^{-D}.
\]
Replacing the even coefficients adds the displayed weighted sum of their
radii.  Finally, each of the six integrals is $1/6$ times its average, so
the six factors $1/6$ sum to one.  Taking a real part cannot increase the
error.  Atoms have no integral tail of this kind.
\end{proof}

Here are explicit rational operations for obtaining the coefficient
enclosures in this lemma.  Write
$\mathcal B(a,r)=\{w\in\mathbb C:|w-a|\le r\}$, where
$a\in\mathbb Q(i)$ and $r\in\mathbb Q_{\ge0}$, and put
$\langle a\rangle_1=|\Re a|+|\Im a|$.  The following enclosure operations
use rational centers and radii:
\begin{align}\label{ind:ball-operations}
 \mathcal B(a,r)+\mathcal B(b,s)&\subseteq\mathcal B(a+b,r+s),\nonumber\\
 \mathcal B(a,r)\mathcal B(b,s)&\subseteq
 \mathcal B\bigl(ab,\langle a\rangle_1s+\langle b\rangle_1r+rs\bigr),\nonumber\\
 \mathcal B(a,r)^{-1}&\subseteq\mathcal B(a^{-1},r/\ell_0^2)
 \quad\text{if }0<\ell_0\le |a|-r.
\end{align}
In the last line $a^{-1}\in\mathbb Q(i)$, and the condition on a rational
$\ell_0$ can be checked by squaring rational quantities.  The bound follows
from $|w^{-1}-a^{-1}|=|w-a|/(|w||a|)$.  These operations allow repeated
use of the same enclosure; possible dependence between its occurrences
does not invalidate inclusion.

Constant exponentials also admit arbitrarily small rational radii.
Use the rational bounds $\mathcal T_n,\mathcal E_n$ from
\eqref{ind:positive-exponential}.  If
$w\in\mathcal B(w_0,r)$, choose rational $q\ge|w_0|$, for example
$q=\langle w_0\rangle_1$, and nonnegative integers $n,s$
with $q<n+2$ and $q+r<s+2$.  Then
\begin{equation}\label{ind:constant-exponential}
 e^w\in\mathcal B\left(\sum_{k=0}^n\frac{w_0^k}{k!},
       \mathcal T_n(q)+r\mathcal E_s(q+r)\right).
\end{equation}
Indeed, the first radius bounds the complex Taylor tail at $w_0$.
The integral first-difference formula for $e^w$ bounds the change from
$w_0$ to $w$ by $r e^{q+r}$, giving the second radius.  Both preconditions
in \eqref{ind:positive-exponential} are explicit; no estimate of a positive
exponential by a truncated sum alone is being used.  For bounded $q$,
$\mathcal T_n(q)\to0$ as $n\to\infty$.

For clarity, the coefficient recurrences needed for the present integrands
are as follows.  If $a(x)=\sum a_kx^k$ and $b(x)=\sum b_kx^k$, multiplication
is convolution,
\begin{equation}\label{ind:convolution}
 (ab)_k=\sum_{v=0}^k a_v b_{k-v}.
\end{equation}
If $a_v\in\mathcal B(\widehat a_v,e_v)$ and
$b_v\in\mathcal B(\widehat b_v,f_v)$, a rational radius for the center
$\sum_{v=0}^k\widehat a_v\widehat b_{k-v}$ is explicitly
\[
 \sum_{v=0}^k\left(\langle\widehat a_v\rangle_1f_{k-v}
       +\langle\widehat b_{k-v}\rangle_1e_v+e_vf_{k-v}\right).
\]
For a reciprocal series $q(x)=a(x)^{-1}$ with $a_0\ne0$, and for an
exponential series $E(x)=e^{g(x)}$, the recurrences are
\begin{align}\label{ind:reciprocal-exponential-recurrences}
 q_0&=a_0^{-1},&
 q_k&=-q_0\sum_{v=1}^k a_vq_{k-v}\quad(k\ge1),\nonumber\\
 E_0&=e^{g_0},&
 E_k&=\frac1k\sum_{v=1}^k v g_v E_{k-v}\quad(k\ge1).
\end{align}
The second line follows by comparing coefficients in $E'=g'E$.
Use \eqref{ind:constant-exponential} for $E_0$ and
\eqref{ind:ball-operations} in these recurrences.  Along with the displayed
convolution radius, this specifies a rational radius for every coefficient,
not merely a formal power series.

More particularly, $d_j=c_j+ih\in\mathbb Q(i)$ and
\begin{equation}\label{ind:kernel-recurrences}
 r_0=d_j^{-1},\quad r_{k+1}=-r_k/d_j,\qquad
 \lambda_k=ir_k/b,\quad z_0=-Br_0-ih,\quad z_k=-Br_k\ (k\ge1)
\end{equation}
are the coefficients of $(t+ih)^{-1}$, $\lambda(t)$, and $z(t)$ at
$t=c_j+x$.  The $r_k,z_k$ are rational complex numbers; an enclosure for
$b^{-1}$ gives the radii for $\lambda_k$.  On piece $j$, put
\[
 L_{n,j}=\sum_{v=j}^6P_v e^{i\pi t_vn},\qquad
 H_{n,j}=\sum_{v=j}^6t_vP_v e^{i\pi t_vn}.
\]
For $\Phi_n(x)=e^{-i\pi n(c_j+x)}$, use
$\phi_{n,0}=e^{-i\pi nc_j}$ and
$\phi_{n,k+1}=(-i\pi n)\phi_{n,k}/(k+1)$.  The density is then the finite
series expression
\begin{equation}\label{ind:density-recurrence}
 f_j(c_j+x)=2\pi\sum_n\Phi_n(x)
   \left\{-\pi c_n\bigl(H_{n,j}-(c_j+x)L_{n,j}\bigr)+id_nL_{n,j}\right\}.
\end{equation}
The constant exponentials in $L_{n,j},H_{n,j},\phi_{n,0}$ are enclosed by
\eqref{ind:constant-exponential}.  For a series $v(x)$ equal to $c_j+x$
or $z(c_j+x)$, the divided powers obey
\begin{equation}\label{ind:divided-power-recurrence}
 V_0(x)=1,\qquad V_{\ell+1}(x)=
       \frac{V_\ell(x)\,i\pi y v(x)}{\ell+1}.
\end{equation}
Finally apply the exponential recurrence to $g=i\pi m(c_j+x)$ or
$g=i\pi m z(c_j+x)$, and multiply the resulting series by the density,
the divided power, and, in the transformed case, $\lambda$.  Truncating
these recurrences after degree $D$ gives every coefficient through $D$
exactly at the level of formal series; omitted higher coefficients cannot
enter a lower convolution coefficient.

The seven possible atoms are handled separately.  At $t_j$, enclose their
masses $e_jCP_j$ and the appropriate factor in
\eqref{cert:analytic-factors} by the same rational operations and constant
exponential bound.  If their respective centers and radii are
$(a_j,r_j)$ and $(g_j,s_j)$, their product has radius
\begin{equation}\label{ind:atom-radii}
 \eta_j=\langle a_j\rangle_1s_j+\langle g_j\rangle_1r_j+r_js_j.
\end{equation}
The sum $\eta=\sum_{j=0}^6\eta_j$ is added to the six-piece error; there is
no Taylor integration tail for an atom.  When $C=0$, all these terms vanish.

All of these radii can be refined effectively.  Alternating rational
arctangent sums in Machin's identity enclose $\pi$ to arbitrarily small
width, and rational bisection of the positive square root of $3/4$ does the
same for $b$.  The kernel denominators have $|c_j+ih|\ge h=2/5$, and
$b>4/5$.  Taking real and imaginary parts of the constant exponential
enclosures also encloses cosine and sine with the same radius.  The nonzero
sine denominators in \eqref{fourier:node-data} have modulus greater than
$1/4$.  Consequently sufficiently refined enclosures certify fixed
positive lower bounds for the moduli of all denominators used above or in
computing $Q_m,D_m$, and separately a positive lower bound for $Q_m$ at
pair division.  For each fixed $D$ there are finitely many operations in
\eqref{ind:convolution}--\eqref{ind:divided-power-recurrence}.  Their
radius formulas tend to zero when their input radii tend to zero, while
\eqref{ind:constant-exponential} has radius tending to zero by refining
its input and increasing its two truncation orders.  Thus one may refine
the rational inputs and exponential orders, recompute the displayed
rational radii, and stop when any prescribed finite collection of positive
coefficient and atom budgets holds.  The convergence just proved ensures
that this rational stopping test terminates.  This is a computability
statement for each specified finite set of inputs, not a claim that any
of its moment sums have been evaluated.

We next make sufficient budgets explicit.  Put
$h_D=\lfloor D/2\rfloor+1$.  For rational $\delta>0$, the conditions
\begin{equation}\label{ind:coefficient-budget}
 e_{j,2r}\le \frac{\delta(2r+1)12^{2r}}{h_D}
 \quad(0\le2r\le D)
\end{equation}
on each piece and moment imply that the weighted coefficient-radius sum in
\eqref{ind:taylor-average-error} is at most $\delta$.

We now combine these integration budgets with the pair estimate
\eqref{ind:pair-stability} and the Bernstein estimate from
Section~\ref{ind:bernstein}.

\begin{proposition}[Two conditional Taylor budgets]\label{ind:taylor-budgets}
The following choices give sufficient rational error enclosures for the
same finite quantities used in Appendix~\ref{sec:certificates}.

For a basis input at $n\in J=\{1,3,4\}$, a row $m\in J$, $y=1$, and
$\ell=0,1$, take $D=63$, $h_D=32$, $M=10^6$, and $\delta=10^{-14}$ in
\eqref{ind:coefficient-budget}.  There are no atoms.  If $Q_m,D_m$ have
rational errors at most $10^{-25}$, either entry of \eqref{cert:pair}
has error less than $2.13\cdot10^{-13}<10^{-12}$.

For the full rational finite-input domain at the start of this subsection,
take $D=255$, $h_D=128$, $M=10^{53}$, and $\delta=4\cdot10^{-15}$ in
\eqref{ind:coefficient-budget}, and require each of the seven atom-product
radii in \eqref{ind:atom-radii} to be at most $4\cdot10^{-15}/7$.
Each $U$ and $X$ moment then has error less than $10^{-14}$.  For node
rows $m\in I_f$, with rational $Q_m,D_m$ errors at most $10^{-25}$, either
entry of \eqref{cert:pair} has error less than $1.8\cdot10^{-14}$.
\end{proposition}

\begin{proof}
The first-block majorant $10^6$ is proved in
Section~\ref{ind:small-kernel}.  Lemma~\ref{ind:taylor-average} gives
\[
 \epsilon\le10^6 2^{-63}+10^{-14}<1.19\cdot10^{-13}.
\]
On the real interval a basis density has modulus below $214$,
$|\lambda|<3$, and $\pi|z|<10$; thus the two exact moments have modulus
below $10^4$.  Substitution of $H=10^4$, $q=7/4$, $d=2.12$, and
$\eta_Q=\eta_D=10^{-25}$ in \eqref{ind:pair-stability} gives
\[
 \frac{(3.12+10^{-25})1.19\cdot10^{-13}+10^4 10^{-25}}{7/4}
 +\frac{3.12\cdot10^4 10^{-25}}{(7/4)^2}<2.13\cdot10^{-13}.
\]
The pointwise majorant $10^{53}$ for the full domain was proved in
Lemma~\ref{cert:moment-error}.  In the second case the total moment radius
is at most
\[
 10^{53}2^{-255}+8\cdot10^{-15}
 <1.728\cdot10^{-24}+8\cdot10^{-15}<10^{-14}.
\]
For the value and first-derivative moments in \eqref{cert:pair}, the real
density is bounded by $1000$, the sum of seven atom moduli by $42$, and
the two factors by $3$ and $30$.  Hence both exact moments are bounded
by $H=4\cdot10^4$.  The right side of \eqref{ind:pair-stability} is then
at most
\[
 \frac{(3.12+10^{-25})10^{-14}+4\cdot10^4 10^{-25}}{7/4}
 +\frac{3.12\cdot4\cdot10^4 10^{-25}}{(7/4)^2}<1.8\cdot10^{-14}.
\]
All inequalities here are comparisons of the displayed rational bounds.
The refinement argument above proves that their coefficient and atom
budgets can be met for any prescribed finite set of these moments.
\end{proof}

For a matrix with $r$ rows and entry error at most $\varepsilon$, the
induced column-norm error is at most $r\varepsilon$, and multiplication
by $W_\eta$ adds the factor $\|W_\eta\|_1$.  In particular $r=6$ for the
$J$ row block and $r=12$ for the rows on $E\cap[7,16]$.  In a tabulated
residual, the other summands in \eqref{cert:residual-bounds} are exact
rationals, so its error is just the corresponding pair error.  Finally,
for the signed half-gaps of Lemma~\ref{signs:bernstein}, the two moment
errors below $10^{-14}$ give a power-coefficient error below
$8\cdot10^{-14}$ and a Bernstein-coefficient error below
$2.32\cdot10^{-12}$, as proved in Section~\ref{ind:bernstein}.
These bounds specify the margins required in a Taylor-based comparison;
they do not themselves assert the matrix, residual, or Bernstein values
printed in Appendix~\ref{sec:certificates}.

\subsection{Direct node-Taylor representation}\label{ind:node-taylor}

There is also a local Taylor representation for the direct cardinal
function that does not integrate its spectral density.  It applies to
the full $\ell^1$ class of \eqref{fourier:cardinal}.  Write
$\mathcal N_0=\{0\}\cup\mathcal N$ and, for $m\in\mathcal N_0$, put
\begin{equation}\label{ind:node-radius}
 \rho_m=\min_{\substack{n\in\mathcal N_0\\n\ne m}}|m-n|.
\end{equation}
The consecutive gaps in one period of $\mathcal N_0$ are
$1,2,1,3,2,3$.  Thus the minimum exists and is positive:
$\rho_m=1$ for $m=0$ and for residues $0,1,3,4$, while $\rho_m=2$
for residues $7,9$.  This is a guaranteed radius to the nearest other
node; it may be smaller than the analytic radius for a particular list
whose corresponding coefficients vanish.

\begin{proposition}\label{ind:node-taylor-proposition}
Let $p=PR$ have a real $\ell^1$ coefficient list as in
\eqref{fourier:list-norm}.  Define
\begin{equation}\label{ind:node-P-coefficients}
 A_k(m)=\frac{P^{(k)}(m)}{k!}
 =\frac1{k!}\sum_{v=-6}^6P_v(i\pi t_v)^k e^{i\pi t_vm}.
\end{equation}
Then $A_0(m)=A_1(m)=0$.  The regular series has coefficients
\begin{equation}\label{ind:node-regular-coefficients}
 r_j(m)=C\mathbf1_{\{j=0\}}+
 (-1)^j\sum_{\substack{n\in\mathcal N_0\\n\ne m}}
 \left(\frac{(j+1)c_n}{(m-n)^{j+2}}
       +\frac{d_n}{(m-n)^{j+1}}\right),
\end{equation}
and converges for $|u|<\rho_m$.  For $0<|u|<\rho_m$, the Laurent
expansion is
\[
 R(m+u)=\frac{c_m}{u^2}+\frac{d_m}{u}
       +\sum_{j=0}^{\infty}r_j(m)u^j.
\]
For every integer $\ell\ge0$ the direct derivative is therefore
\begin{equation}\label{ind:node-derivatives}
 \frac{p^{(\ell)}(m)}{\ell!}
 =c_mA_{\ell+2}(m)+d_mA_{\ell+1}(m)
       +\sum_{k=2}^{\ell}A_k(m)r_{\ell-k}(m),
\end{equation}
where the last sum is empty for $\ell=0,1$.
\end{proposition}

\begin{proof}
For $n\ne m$, the geometric series for $(m-n+u)^{-1}$ and its derivative
give
\[
 \frac1{m-n+u}=\sum_{j\ge0}\frac{(-1)^ju^j}{(m-n)^{j+1}},\qquad
 \frac1{(m-n+u)^2}=\sum_{j\ge0}
       \frac{(-1)^j(j+1)u^j}{(m-n)^{j+2}},
\]
whenever $|u|<|m-n|$.  To justify summing these series over the infinite
list, fix $0\le r<\rho_m$ and set $q=r/\rho_m<1$.  Since
$|m-n|\ge\rho_m$ for every other node, the sum of absolute values on
$|u|\le r$ is at most
\[
 \frac{\sum_{n\ne m}|c_n|}{\rho_m^2}
       \sum_{j\ge0}(j+1)q^j
 +\frac{\sum_{n\ne m}|d_n|}{\rho_m}\sum_{j\ge0}q^j
 =
 \frac{\sum_{n\ne m}|c_n|}{\rho_m^2(1-q)^2}
 +\frac{\sum_{n\ne m}|d_n|}{\rho_m(1-q)}<\infty.
\]
This proves absolute uniform convergence on every smaller closed disk,
permits exchanging the node and Taylor sums, and gives
\eqref{ind:node-regular-coefficients}.  It also proves precisely why the
nearest-other-node distance is a valid common radius.

The finite Fourier expansion \eqref{fourier:polynomial} gives
\eqref{ind:node-P-coefficients}, and the double zero of $P$ gives
$A_0=A_1=0$.  Multiply its entire Taylor series by the convergent Laurent
series for $R$.  The two principal terms are canceled by the double zero;
collecting the coefficient of $u^\ell$ gives
\eqref{ind:node-derivatives}.  In particular $A_2=Q_m$ and $A_3=Q_mD_m$,
so the cases $\ell=0,1$ reproduce \eqref{fourier:node-values}.
\end{proof}

For finite rational input, the constants in this direct formula can be
enclosed by the rational operations above.  For an arbitrary $\ell^1$
input the formula is an exact convergence statement; a computable error
for truncating its node sum additionally requires a bound for that list's
tail.  It represents the direct function $p$ only.  The transformed
function $K_p$ still has the kernel representation
\eqref{fourier:kernel}, and this local identity neither places $K_p$ in
the cardinal class nor replaces the exact-derivative certificates for the
tabulated functions.
